\begin{figure}[H]
  \centering
  \begin{tikzpicture}[
      cell/.style={draw=codeviolet, fill=codeback, rounded corners=1pt,
        minimum width=1.22cm, minimum height=0.70cm, inner sep=1pt,
        font=\ttfamily\footnotesize, anchor=center},
      title/.style={font=\small\bfseries, anchor=center},
      note/.style={font=\footnotesize, anchor=center},
      rowlabel/.style={font=\ttfamily\footnotesize, anchor=east},
      valid/.style={circle, draw=codeviolet, fill=codeviolet,
        minimum size=2.4pt, inner sep=0pt},
      invalid/.style={circle, draw=codeviolet, fill=white,
        minimum size=2.4pt, inner sep=0pt},
      arrow/.style={-{Latex[length=2mm]}, thick, draw=codeviolet}
    ]
    % (a) IMEM: palabras de 32 bits, direccionadas por PC >> 2.
    \node[title] at (2.05,1.45) {(a) IMEM};
    \node[note] at (2.05,1.05) {256 palabras de 32 bits};
    \foreach \i/\label in {0/+0, 1/+1, 2/+2, 3/+3} {
      \node[note] at ({0.5+\i*1.3},0.52) {\label};
      \node[cell] (im\i) at ({0.5+\i*1.3},0) {b\i};
    }
    \node[note, anchor=east] at (-0.15,0) {PC};
    \node[note] at (2.05,-0.90) {palabra de 32 bits};
    \node[note] at (2.05,-1.30) {índice de palabra = \code{PC >> 2}};

    % (b) DMEM: cuatro lanes de bytes y una marca de validez por byte.
    \node[title] at (9.95,1.45) {(b) DMEM};
    \node[note] at (9.95,1.05) {4 lanes de bytes $\times$ 512 filas};
    \foreach \lane/\label in {0/0, 1/1, 2/2, 3/3} {
      \node[note] at ({8.15+\lane*1.30},0.52) {lane \label};
    }
    \foreach \row/\addr/\ba/\bb/\bc/\bd/\validflag in {
        0/0x0000/00/00/00/00/0,
        1/0x0004/78/56/34/12/1,
        2/0x0008/00/00/00/00/0,
        3/0x000C/00/00/00/00/0
      } {
        \node[rowlabel] at (7.55,{-\row*1.05}) {\addr};
        \foreach \lane/\cellvalue in {0/\ba, 1/\bb, 2/\bc, 3/\bd} {
          \node[cell] at ({8.15+\lane*1.30},{-\row*1.05}) {\cellvalue};
          \node[invalid] at ({7.75+\lane*1.30},{-\row*1.05+0.23}) {};
        }
      }
    \foreach \lane in {0, 1, 2, 3} {
      \node[valid] at ({7.75+\lane*1.30},{-0.82}) {};
    }
    \draw[codeviolet, thick, rounded corners=2pt] (7.45,-1.50) rectangle (12.75,-0.60);
    \node[note, anchor=west] at (13.15,-1.05) {$A=0x0004$};
  \end{tikzpicture}
  \caption{Organización de las memorias. Cada celda representa un byte; en
  DMEM, el círculo superior izquierdo indica la validez de ese byte: lleno si
  fue escrito y vacío si todavía se lee como cero. La fila resaltada muestra la
  palabra \code{0x12345678} almacenada en \code{0x0004}.}
  \label{fig:memoria-organizacion}
\end{figure}
